(1) 令 \(x=pE/(k_BT)\)。对球面立体角积分 \(z_{\rm or}=\int_0^{2\pi}d\varphi\int_0^\pi e^{x\cos\theta}\sin\theta d\theta=4\pi\sinh x/x\)，取 \(x=0\) 的连续值为 \(4\pi\)。\(F_{\rm or}=-Nk_BT\ln z_{\rm or}\)。

(2) \(\langle\cos\theta\rangle=\partial_x\ln z_{\rm or}=\coth x-1/x\equiv\mathcal L(x)\)，总极化 \(P_{\parallel}=Np\mathcal L(x)\)。弱场 \(\mathcal L(x)=x/3+O(x^3)\)，故 \(\partial P/\partial E|_{E=0}=Np^2/(3k_BT)\)；强场 \(x\gg1\) 时 \(\mathcal L\to1\)，总极化趋 \(Np\)。

(3) \(U_{\rm or}=-NpE\mathcal L(x)=-Nk_BTx\mathcal L(x)\)。固定 \(E\) 时 \(dx/dT=-x/T\)，所以 \(C_{E}=Nk_Bx^2\mathcal L'(x)=Nk_B[1-x^2/\sinh^2x]\ge0\)。弱场 \(C_E\simeq Nk_Bx^2/3\)；强场 \(C_E\to Nk_B\)，对应两个小角度自由度。

(4) \(S_{\rm or}=(U-F)/T=Nk_B[\ln(4\pi\sinh x/x)-x\mathcal L(x)]\)。求导 \(\partial_xS=-Nk_Bx\mathcal L'(x)\le0\)，故固定 \(T\) 增大场会降低取向熵。\(x\to0\) 时 \(S\to Nk_B\ln4\pi\)，是每一自由偶极的均匀球面取向。
